\subsection{Inflation on standard affinoids}\label{subsec:fc-line-inflation}

\begin{definition}[Inflation with a fixed centre]\label{def:fc-centered-inflation}
\uses{conv:fc-line, analytic:actual-polynomial-gauss-inflation}
\lean{EnvelopeProperty.Berkovich.RelativeSpectrum.polynomialTranslation, EnvelopeProperty.Berkovich.RelativeSpectrum.fcCenteredInflation}
\leanok
Let $v$ be a real-valued ultrametric valuation on a commutative ring $F$, and
let $x$ be a multiplicative seminorm on $F[T]$ extending $v$.
For $a\in F$ put $t_a(x)(P)=x(P(T+a))$. Define
$H_a(x,s)=t_a(H(t_{-a}(x),s))$, where
$H(x,s)(P)$ is the Gauss seminorm of $P(T+S)$ in $S$, with coefficient
seminorm $x$ and radius $s$.
\end{definition}

\begin{lemma}[Exact linear values and endpoints]\label{lem:fc-inflation-linear}
\uses{def:fc-centered-inflation}
\lean{EnvelopeProperty.Berkovich.RelativeSpectrum.continuous_fcCenteredInflation_prod, EnvelopeProperty.Berkovich.RelativeSpectrum.fcCenteredInflation_val_X_sub_C, EnvelopeProperty.Berkovich.RelativeSpectrum.fcCenteredInflation_zero, EnvelopeProperty.Berkovich.RelativeSpectrum.fcCenteredInflation_eq_gauss}
\leanok
The map $(x,s)\mapsto H_a(x,s)$ is continuous, $H_a(x,0)=x$, and
\[
 |T-b|_{H_a(x,s)}=\max\{|T-b|_x,s\}.
\]
If $|T-a|_x\le s$, then $H_a(x,s)=\zeta_{a,s}$.
\end{lemma}

\begin{proof}
Polynomial translation is an algebra automorphism, and every value of a Gauss
seminorm is the maximum of finitely many continuous coefficient values times
powers of the radius. This proves continuity. At radius zero, taking the
constant coefficient of $P(T+S)$ gives $P(T)$; the two translations cancel.
For $P=T-b$ the translated polynomial in $S$ has constant coefficient $T-b$
and coefficient $1$ in degree one, giving the displayed maximum.
If $|T-a|_x\le s$, translate $a$ to zero. Translation by the coefficient
element $T$ fixes the radius-$s$ Gauss seminorm because its coefficient
seminorm is at most $s$. Restriction back to $F[T]$ is therefore the
radius-$s$ Gauss seminorm on $F[T]$. Translating back gives $\zeta_{a,s}$.
\end{proof}

\begin{lemma}[Linear values determine points]\label{lem:fc-linear-determination}
\uses{conv:fc-line, analytic:actual-algclosed-disc-branch-cover}
\lean{EnvelopeProperty.Berkovich.RelativeSpectrum.fc_ext_linear}
\leanok
Over an algebraically closed field, two points of the analytic affine line
with the same values on every $T-a$ are equal.
\end{lemma}

\begin{proof}
Every polynomial is its leading coefficient times the product of its linear
factors, with roots repeated according to multiplicity. The seminorms have
the same values on constants because both extend the prescribed field
valuation. Multiplicativity therefore gives equality on every polynomial.
\end{proof}

\begin{lemma}[Topology from linear coordinates]\label{lem:fc-linear-topology}
\uses{lem:fc-linear-determination}
\lean{EnvelopeProperty.Berkovich.RelativeSpectrum.fc_continuous_iff_linear, EnvelopeProperty.Berkovich.RelativeSpectrum.fc_continuous_centeredGauss_prod}
\leanok
Over an algebraically closed valued field, a map to the analytic affine line
is continuous if and only if its evaluations on every $T-a$ are continuous.
If the valuation is the norm of the field, $(a,r)\mapsto\zeta_{a,r}$ is jointly
continuous.
\end{lemma}

\begin{proof}
Continuity of every polynomial value implies continuity of the linear values.
Conversely, factor each polynomial into its leading coefficient and finitely
many linear factors. Its value is a fixed constant times a finite product of
linear values, hence is continuous. For the Gauss points each linear value is
$\max\{|a-b|,r\}$, which is continuous in $(a,r)$.
\end{proof}

\begin{lemma}[Strict residue discs]\label{lem:fc-strict-residue-disc}
\uses{lem:fc-linear-determination, lem:fc-inflation-linear, analytic:actual-algclosed-disc-branch-cover}
\lean{EnvelopeProperty.Berkovich.RelativeSpectrum.fc_exists_strict_residue_disc}
\leanok
If $x\in D(a,r)$ and $x\ne\zeta_{a,r}$, there is $b\in F$ with
$|a-b|\le r$ and $|T-b|_x<r$.
\end{lemma}

\begin{proof}
Suppose no such $b$ exists. For every $b$ the ultrametric inequality gives
$|T-b|_x\le\max\{|T-a|_x,|a-b|\}\le\max\{r,|a-b|\}$.
If $|a-b|\le r$, the supposition gives $|T-b|_x\ge r$.
If $|a-b|>r$, write $a-b=(T-b)-(T-a)$: its value is at most
$\max\{|T-b|_x,|T-a|_x\}$, and the second entry is strictly smaller
than $|a-b|$, so $|T-b|_x\ge|a-b|$. Thus every $|T-b|_x$ equals
$\max\{r,|a-b|\}=|T-b|_{\zeta_{a,r}}$. Linear determination gives
$x=\zeta_{a,r}$, a contradiction.
\end{proof}

\begin{lemma}[Interior of a closed disc]\label{lem:fc-closed-disc-interior}
\uses{lem:fc-strict-residue-disc}
\lean{EnvelopeProperty.Berkovich.RelativeSpectrum.fc_interior_closedDisc}
\leanok
The topological interior of $D(a,r)$ is $D(a,r)\setminus\{\zeta_{a,r}\}$.
\end{lemma}

\begin{proof}
For a point $x\ne\zeta_{a,r}$ in the closed disc, the preceding lemma gives
$b$ such that $x\in D^-(b,r)$ and $|a-b|\le r$. For every $y$ in this open
disc, the ultrametric inequality gives
$|T-a|_y\le\max\{|T-b|_y,|b-a|\}\le r$. Thus $x$ is interior.
On the other hand, continuity of $s\mapsto\zeta_{a,s}$ implies that every
neighbourhood of $\zeta_{a,r}$ contains $\zeta_{a,s}$ for some $s>r$.
That point lies outside $D(a,r)$, so the Gauss point is not interior.
\end{proof}

\begin{lemma}[Interior of an open-disc complement]\label{lem:fc-disc-exterior-interior}
\uses{lem:fc-strict-residue-disc}
\lean{EnvelopeProperty.Berkovich.RelativeSpectrum.fc_interior_discExterior}
\leanok
For $r>0$ the topological interior of $\{|T-a|\ge r\}$ is that set with
$\zeta_{a,r}$ removed.
\end{lemma}

\begin{proof}
Every point with $|T-a|>r$ belongs to the open set defined by that strict
inequality. At a point $x\ne\zeta_{a,r}$ with $|T-a|_x=r$, choose a strict
residue disc $D^-(b,r)$ as above. The inequalities
$r=|T-a|_x\le\max\{|T-b|_x,|b-a|\}$ and $|T-b|_x<r$ imply
$|a-b|\ge r$. For any $y\in D^-(b,r)$,
$r\le|a-b|\le\max\{|T-a|_y,|T-b|_y\}$, so $|T-a|_y\ge r$.
The residue disc therefore supplies an open neighbourhood in the complement.
The Gauss point is approached by $\zeta_{a,s}$ with $s<r$; these lie in the
removed open disc, proving that it is not interior.
\end{proof}

\begin{remark}[Degenerate radii and repeated boundary points]\label{rmk:fc-standard-degenerate}
\uses{def:fc-standard-affinoid}
\lean{EnvelopeProperty.Berkovich.RelativeSpectrum.fc_closedDisc_zero, EnvelopeProperty.Berkovich.RelativeSpectrum.FCStandardAffinoid.no_holes_of_radius_zero, EnvelopeProperty.Berkovich.RelativeSpectrum.FCStandardAffinoid.holeGauss_eq_outer_of_radius_eq}
\leanok
A zero-radius outer disc has no positive-radius holes and consists of its
evaluation point. If a hole has the same radius as the outer disc, its Gauss
point equals the outer Gauss point. Boundary points are a finite set, so
coincidences are retained only once. The formulae for compactness, paths and
interiors also apply when there are no holes.
\end{remark}

\begin{lemma}[Compactness and paths in standard affinoids]\label{lem:fc-standard-inflation}
\uses{def:fc-standard-affinoid, lem:fc-inflation-linear}
\lean{EnvelopeProperty.Berkovich.RelativeSpectrum.FCStandardAffinoid.isCompact_carrier, EnvelopeProperty.Berkovich.RelativeSpectrum.FCStandardAffinoid.isPathConnected_carrier}
\leanok
Every standard affinoid is compact and path connected. For $x\in U$ and
$0\le s\le r_0$, the point $H_{a_0}(x,s)$ lies in $U$, and the resulting
path joins $x$ to $\zeta_{a_0,r_0}$.
\end{lemma}

\begin{proof}
The closed disc $D(a_0,r_0)$ is the image under polynomial translation of
the compact coordinate disc $D(0,r_0)$. Each condition
$|T-a_j|\ge r_j$ is closed, so their intersection with that disc is compact.
Under inflation the outer coordinate value is
$\max\{|T-a_0|_x,s\}\le r_0$, and each hole coordinate value is
$\max\{|T-a_j|_x,s\}\ge r_j$. Thus inflation stays in $U$.
At radius zero it equals $x$ and at radius $r_0$ it equals the outer Gauss
point. Parametrizing $s=t r_0$ for $0\le t\le1$ gives a continuous path.
Reversing one such path and following another joins any two points of $U$.
\end{proof}

\begin{lemma}[The exact finite boundary]\label{lem:fc-standard-boundary}
\uses{def:fc-standard-affinoid, lem:fc-closed-disc-interior, lem:fc-disc-exterior-interior}
\lean{EnvelopeProperty.Berkovich.RelativeSpectrum.FCStandardAffinoid.boundary_subset, EnvelopeProperty.Berkovich.RelativeSpectrum.FCStandardAffinoid.interior_carrier, EnvelopeProperty.Berkovich.RelativeSpectrum.FCStandardAffinoid.frontier_carrier}
\leanok
The specified finite Gauss boundary of a standard affinoid lies in the affinoid
and equals its topological frontier. Its topological interior is
$U\setminus\partial U$.
\end{lemma}

\begin{proof}
The outer Gauss point belongs to every hole complement since its value on
$T-a_j$ is at least $r_0\ge r_j$. The Gauss point of hole $j$ belongs to the
outer disc by the centre and radius bounds. If it were in another open hole,
then the rigid point $a_j$ would also be in that hole: its value on each
$T-a_i$ is no greater than the value at $\zeta_{a_j,r_j}$. But $a_j$ belongs
to its own open hole, contradicting disjointness. Thus all boundary points
belong to $U$.
Interior commutes with finite intersections. Apply the two preceding
interior formulae to the outer disc and the finitely many hole complements;
the result removes precisely their Gauss points. The set $U$ is closed because
it is a finite intersection of closed coordinate inequalities, so its frontier
is $U\setminus\operatorname{int}U=\partial U$.
\end{proof}

\begin{lemma}[Interior components approach the boundary]\label{lem:fc-standard-component-boundary}
\uses{lem:fc-standard-inflation, lem:fc-standard-boundary}
\lean{EnvelopeProperty.Berkovich.RelativeSpectrum.FCStandardAffinoid.exists_boundary_mem_closure_component}
\leanok
For every $x\in U\setminus\partial U$, the closure of its connected
component in $U\setminus\partial U$ contains a point of $\partial U$.
\end{lemma}

\begin{proof}
Let $p\colon[0,1]\to U$ be the inflation path from $x$ to the outer Gauss
point. The finite boundary is closed. Its inverse image under $p$ is a
nonempty closed subset of the compact interval, since it contains $1$.
It therefore has a least element $t$. Since $p(0)=x$ is not in the boundary,
$t>0$. The connected set $p([0,t))$ contains $x$ and avoids the boundary,
so it lies in the connected component of $x$ in $U\setminus\partial U$.
Continuity gives $p(t)\in\overline{p([0,t))}$. This point belongs to the
boundary and to the closure of the required component.
\end{proof}

\begin{corollary}[Paths in the entire analytic line]\label{cor:fc-line-paths}
\uses{analytic:valued-ring-affine-line-contractible}
\lean{EnvelopeProperty.Berkovich.RelativeSpectrum.fc_joinedIn_line, EnvelopeProperty.Berkovich.RelativeSpectrum.fc_isPathConnected_line, EnvelopeProperty.Berkovich.RelativeSpectrum.fc_isConnected_line}
\leanok
The analytic affine line is path connected, hence connected.
\end{corollary}

\begin{proof}
The whole polynomial inflation homotopy deforms a point $x$ to the Gauss
point of radius $|T|_x$ with the same coefficient restriction. Shrinking that
radius to zero then gives a contraction to the evaluation point at zero,
because the coefficient restriction is the prescribed valuation. Following
this contraction and reversing it joins any two points.
\end{proof}

\begin{lemma}[Retraction to the finite radial skeleton]\label{lem:fc-skeleton-retraction}
\lean{EnvelopeProperty.Berkovich.RelativeSpectrum.FCStandardAffinoid.radialRetraction, EnvelopeProperty.Berkovich.RelativeSpectrum.FCStandardAffinoid.radialRetraction_eq_self}
\leanok
\uses{lem:fc-inflation-linear, lem:fc-linear-determination, lem:fc-standard-boundary}
The radial description of the skeleton is
\[
 \Sigma(U)=\{\zeta_{a_0,r_0}\}\cup
 \bigcup_{j=1}^m\{\zeta_{a_j,s}:r_j\le s\le r_0\}.
\]
Set $s_U(x)=\min\bigl(\{r_0\}\cup\{|T-a_j|_x:1\le j\le m\}\bigr)$.
The map $\tau_U(x)=H_{a_0}(x,s_U(x))$ is a continuous retraction
$U\to\Sigma(U)$.
\end{lemma}

\begin{proof}
The exact linear-value formula shows that $H_a(x,s)$ is independent of $a$:
two choices have the same values on every linear polynomial and hence are
equal. The radius $s_U$ is continuous, as the minimum of finitely many
continuous functions. Thus $\tau_U$ is continuous and stays in $U$.
The finite minimum is attained. If it equals the outer radius, the endpoint
formula gives $\tau_U(x)=\zeta_{a_0,r_0}$. Otherwise choose $j$ with
$s_U(x)=|T-a_j|_x$. This radius lies in $[r_j,r_0]$ because $x\in U$;
centre independence and the endpoint formula give
$\tau_U(x)=\zeta_{a_j,s_U(x)}$. This proves that its image lies in $\Sigma(U)$.
Every displayed radial segment lies in $U$: the outer coordinate is at most
$r_0$, and its value on every hole coordinate is no smaller than the value
at $\zeta_{a_j,r_j}$, which belongs to $U$ by the boundary lemma.
At a point $\zeta_{a_j,s}$ of such a segment every hole-coordinate value is
at least $s$, while the $j$th equals $s$ and $r_0\ge s$. Hence
$s_U(\zeta_{a_j,s})=s$. Inflating that Gauss point by radius $s$ fixes it,
so $\tau_U$ fixes the segment. The same argument at the outer Gauss point
uses that all its hole-coordinate values equal $r_0$; the minimum is $r_0$,
and the endpoint formula fixes it. Therefore $\tau_U$ fixes all of $\Sigma(U)$.
\end{proof}
